The massive growth of Machine Learning in recent years has led to an unprecedented movement of data to centralized servers. Beyond creating severe bandwidth bottlenecks, this centralized approach raises significant privacy and regulatory concerns, as AI models frequently learn from highly sensitive or personal data. To address these challenges, Federated Learning (FL) \cite{mcmahan2017communication} was introduced. The fundamental premise of FL is to execute model training locally where the data resides, thereby eliminating the need to transmit raw datasets to a central repository. In this learning paradigm, a central server initializes a global model and distributes it to participating clients. Each client trains the model locally using its private dataset and then sends only the updated parameters back to the server for aggregation. This paradigm significantly reduces bandwidth consumption while preserving data privacy.

Traditional FL effectively solves the data privacy dilemma, but its dependence on a central server introduces a Single Point of Failure (SPOF) and creates new communication bottlenecks as networks scale. To mitigate these structural weaknesses, the paradigm is evolving towards Decentralized Federated Learning (DFL) \cite{hegedus2019gossip,vanhaesebrouck2017decentralized}, often referred to as Gossip Learning (GL) \cite{ormandi2013gossip}. In DFL, the central orchestrator is eliminated entirely. Instead, clients communicate directly with their topological neighbors in a Peer-to-Peer (P2P) network, continuously exchanging and averaging model parameters to achieve a global consensus (see Figure \ref{fig:DFL_consensus} for a visual representation).

\begin{figure}[t]
    \centering
    \subfloat[Centralized Federated Learning\label{fig:fl_sub}]{%
        \includegraphics[width=0.45\textwidth]{figures/fl_graph.png}%
    }\hfill
    \subfloat[Decentralized Federated Learning\label{fig:dfl_sub}]{%
        \includegraphics[width=0.45\textwidth]{figures/dfl_graph.png}%
    }
    \caption{Structural comparison between the centralized orchestrator dependency in standard FL and the peer-to-peer topology of DFL.}
    \label{fig:DFL_consensus}
\end{figure}

The removal of a central server, while improving fault tolerance and scalability, shifts the security issues primarily onto the network's topology. In centralized FL, the server can act as a robust gatekeeper, using the mass consensus of all the clients to dilute malicious updates. In DFL, the propagation of a poisoned model or a backdoor cannot be controlled or mitigated that easily, leaving the structural layout of the network as one of the unique factors influencing security.

This structural shift exposes a critical research gap regarding how different network topologies influence the viability of localized poisoning attacks. In decentralized environments, an adversary attempting to compromise the performance of its peers faces a fundamental "Effectiveness vs. Stealthiness Trade-off" \cite{yang2025stealthy, syros2024backdoor}. To successfully spread a backdoor and overcome the natural dilution of network aggregation, an attacker must apply an aggressive scaling factor to their malicious updates \cite{bagdasaryan2020how}. However, if the attack is disproportionately aggressive relative to the network's structural resistance, the attacker's local neural network suffers from Catastrophic Forgetting \cite{french1999catastrophic}. This destroys the node's legitimate performance (Clean Accuracy), exposing the attacker and being easily detectable by standard anomaly detection systems.

This thesis aims to empirically demonstrate how network topology dictates the success rate of backdoor poisoning attacks in decentralized machine learning environments. The specific objectives of this research are:

\begin{itemize}
    \item To evaluate how vulnerable different network topologies are to continuous backdoor attacks.
    \item To establish a mathematically controlled adversarial baseline by stabilizing the effectiveness versus stealthiness trade-off, ensuring that network topology remains the sole independent variable governing backdoor propagation.
    \item To investigate how the attacker's position affects the spread of the infection, specifically comparing an attacker at a network bottleneck (Gateway Attacker) versus one isolated inside a community (Peripheral Attacker) \cite{piaseczny2025adversarial}.
    \item To evaluate the balance between Collaboration Gain and Security Risk, investigating whether specific network architectures (such as community models) can naturally mitigate attacks without significantly degrading the overall learning process.
    \item To comprehensively validate these findings across different levels of data complexity, first utilizing a controlled academic dataset (MNIST) \cite{lecun1998gradient} to isolate the topological mechanics, and subsequently deploying a real-world industrial dataset (NEU-DET) \cite{bao2019neudet} to prove that these vulnerabilities persist in practical applications.
\end{itemize}

The rest of this document is organized as follows: Chapter \ref{ch:theoretical framework} establishes the theoretical framework, detailing the transition from FL to DFL. Chapter \ref{ch:threat landscape} explores the threat landscape in decentralized environments, outlining specific vulnerabilities and examining the fundamental mechanics of backdoor poisoning attacks. Chapter \ref{ch:methodology} outlines the methodology, defining the threat model and the novel topological heuristics used for the backdoor implementation. Chapter \ref{ch:experimentation} details the experimental setup, including the evaluated network topologies, dataset, and simulation constraints. Chapter \ref{ch:results} presents the experimental results, providing a comparative analysis of the attacker's performance and stealthiness across different architectural constraints. Finally, Chapter \ref{ch:conclusions} summarizes the core findings and proposes future lines of research in decentralized network security.